% GENERATED FILE. Edit canonical lessons, then run npm run generate:books.
% canonical-source-hash: fnv1a64:fecc8a0ee31e3d21
% canonical-lessons: RU-C514-katok, RU-C514-figurist, RU-C514-skolzit, RU-C514-lasty, RU-C514-pleskatsya

\chapter{On the Ice and in the Water}
\label{ch:ru-on-the-ice-and-in-the-water}
\hlchaptermodality{514}

\begin{chapteropening}
\textbf{By the end of this chapter:} \emph{I can talk about skating and splashing about in the water.}

Complete the last lesson of chapter 514: I can talk about skating and splashing about in the water.
\end{chapteropening}

\section[\texorpdfstring{\ru{каток} (katók)}{каток (katók)}]{\ru{каток} (katók) — a skating rink}

\label{lesson:RU-C514-katok}



\begin{quote}
Before the new one: say the Russian for dumbbells, then the Russian for stamina.
\end{quote}

\subsection*{You'll want to know: \ru{каток}}
\textbf{\ru{каток}} (\emph{katók}) — \textquotedblleft{}a skating rink\textquotedblright{}.

\begin{quote}
\ru{Зимой} \ru{мы} \ru{часто} \ru{ходим} \ru{на} \ru{каток}.
\end{quote}

\begin{quote}
\emph{Zimóy my chásto khódim na katók.}
\end{quote}

\begin{quote}
In winter we often go to the skating rink.
\end{quote}

Masculine; genitive \ru{катка}.

\begin{cousinweb}
\emph{\ru{каток}} is built on the root of \emph{\ru{катать}} 'to roll' and \emph{\ru{кататься}} 'to ride, to skate'; EtymDB has no entry.
\end{cousinweb}

\begin{grammarlens}[title={What you've built}]
One more word for telling, explaining and describing at B1.
\end{grammarlens}

\subsection*{Your turn}
\begin{itemize}
\raggedright
  \item \emph{Say it:} \emph{katók}
  \item \emph{Say it:} \emph{katók}, once more
  \item \emph{Say it:} a sentence of your own with \emph{katók}
  \item \emph{Recall:} say the Russian for dumbbells, then the Russian for stamina, then say \emph{katók} again
\end{itemize}

\begin{tcolorbox}[breakable,colback=teal!4,colframe=teal!35!black,title={Before you move on}]
What does \ru{каток} mean? (\textquotedblleft{}A skating rink\textquotedblright{}.) Say it once more.
\end{tcolorbox}
\section[\texorpdfstring{\ru{фигурист} (figuríst)}{фигурист (figuríst)}]{\ru{фигурист} (figuríst) — a figure skater}

\label{lesson:RU-C514-figurist}



\begin{quote}
Before the new one: say the Russian for stamina, then the Russian for a skating rink.
\end{quote}

\subsection*{You'll want to know: \ru{фигурист}}
\textbf{\ru{фигурист}} (\emph{figuríst}) — \textquotedblleft{}a figure skater\textquotedblright{}.

\begin{quote}
\ru{Фигурист} \ru{выступил} \ru{очень} \ru{красиво}.
\end{quote}

\begin{quote}
\emph{Figuríst výstupil óchen krasívo.}
\end{quote}

\begin{quote}
The figure skater performed beautifully.
\end{quote}

Masculine; feminine \ru{фигуристка}.

\begin{cousinweb}
\emph{\ru{фигурист}} is built on \emph{\ru{фигура}} 'figure', usually traced to Latin \emph{figūra} 'a shape'; EtymDB has no entry.
\end{cousinweb}

\begin{grammarlens}[title={What you've built}]
One more word for telling, explaining and describing at B1.
\end{grammarlens}

\subsection*{Your turn}
\begin{itemize}
\raggedright
  \item \emph{Say it:} \emph{figuríst}
  \item \emph{Say it:} \emph{figuríst}, once more
  \item \emph{Say it:} a sentence of your own with \emph{figuríst}
  \item \emph{Recall:} say the Russian for stamina, then the Russian for a skating rink, then say \emph{figuríst} again
\end{itemize}

\begin{tcolorbox}[breakable,colback=teal!4,colframe=teal!35!black,title={Before you move on}]
What does \ru{фигурист} mean? (\textquotedblleft{}A figure skater\textquotedblright{}.) Say it once more.
\end{tcolorbox}
\section[\texorpdfstring{\ru{скользить} (skolzít)}{скользить (skolzít)}]{\ru{скользить} (skolzít) — to slide, to glide}

\label{lesson:RU-C514-skolzit}



\begin{quote}
Before the new one: say the Russian for a skating rink, then the Russian for a figure skater.
\end{quote}

\subsection*{You'll want to know: \ru{скользить}}
\textbf{\ru{скользить}} (\emph{skolzít}) — \textquotedblleft{}to slide, to glide\textquotedblright{}.

\begin{quote}
\ru{Фигуристы} \ru{легко} \ru{скользят} \ru{по} \ru{льду}.
\end{quote}

\begin{quote}
\emph{Figurísty legkó skolzyát po ldu.}
\end{quote}

\begin{quote}
The figure skaters glide easily over the ice.
\end{quote}

Imperfective; present \ru{скольжу}, \ru{скользишь} (\ru{з} becomes \ru{ж}). Perfective \ru{скользнуть}.

\begin{cousinweb}
EtymDB lists \emph{\ru{скользить}} without a source; it is usually traced to Proto-Slavic \emph{skol\ru{ь}zěti}.
\end{cousinweb}

\begin{grammarlens}[title={What you've built}]
One more word for telling, explaining and describing at B1.
\end{grammarlens}

\subsection*{Your turn}
\begin{itemize}
\raggedright
  \item \emph{Say it:} \emph{skolzít}
  \item \emph{Say it:} \emph{skolzít}, once more
  \item \emph{Say it:} a sentence of your own with \emph{skolzít}
  \item \emph{Recall:} say the Russian for a skating rink, then the Russian for a figure skater, then say \emph{skolzít} again
\end{itemize}

\begin{tcolorbox}[breakable,colback=teal!4,colframe=teal!35!black,title={Before you move on}]
What does \ru{скользить} mean? (\textquotedblleft{}To slide, to glide\textquotedblright{}.) Say it once more.
\end{tcolorbox}
\section[\texorpdfstring{\ru{ласты} (lásty)}{ласты (lásty)}]{\ru{ласты} (lásty) — flippers}

\label{lesson:RU-C514-lasty}



\begin{quote}
Before the new one: say the Russian for a figure skater, then the Russian for to slide.
\end{quote}

\subsection*{You'll want to know: \ru{ласты}}
\textbf{\ru{ласты}} (\emph{lásty}) — \textquotedblleft{}flippers\textquotedblright{}.

\begin{quote}
\ru{В} \ru{ластах} \ru{я} \ru{плаваю} \ru{гораздо} \ru{быстрее}.
\end{quote}

\begin{quote}
\emph{V lástakh ya plávayu gorázdo bystréye.}
\end{quote}

\begin{quote}
I swim much faster in flippers.
\end{quote}

Plural of \emph{\ru{ласт}}.

\begin{cousinweb}
EtymDB has no entry for \emph{\ru{ласты}}, and the origin is not settled.
\end{cousinweb}

\begin{grammarlens}[title={What you've built}]
One more word for telling, explaining and describing at B1.
\end{grammarlens}

\subsection*{Your turn}
\begin{itemize}
\raggedright
  \item \emph{Say it:} \emph{lásty}
  \item \emph{Say it:} \emph{lásty}, once more
  \item \emph{Say it:} a sentence of your own with \emph{lásty}
  \item \emph{Recall:} say the Russian for a figure skater, then the Russian for to slide, then say \emph{lásty} again
\end{itemize}

\begin{tcolorbox}[breakable,colback=teal!4,colframe=teal!35!black,title={Before you move on}]
What does \ru{ласты} mean? (\textquotedblleft{}Flippers\textquotedblright{}.) Say it once more.
\end{tcolorbox}
\section[\texorpdfstring{\ru{плескаться} (pleskátsya)}{плескаться (pleskátsya)}]{\ru{плескаться} (pleskátsya) — to splash about}

\label{lesson:RU-C514-pleskatsya}



\begin{quote}
Before the new one: say the Russian for to slide, then the Russian for flippers.
\end{quote}

\subsection*{You'll want to know: \ru{плескаться}}
\textbf{\ru{плескаться}} (\emph{pleskátsya}) — \textquotedblleft{}to splash about\textquotedblright{}.

\begin{quote}
\ru{Дети} \ru{весь} \ru{день} \ru{плескались} \ru{в} \ru{реке}.
\end{quote}

\begin{quote}
\emph{Déti ves den pleskális v reké.}
\end{quote}

\begin{quote}
The children splashed about in the river all day.
\end{quote}

Imperfective; present \ru{плещусь}, \ru{плещешься} (\ru{ск} becomes \ru{щ}).

\begin{cousinweb}
\emph{\ru{плескаться}} is built on \emph{\ru{плеск}} 'a splash', an imitative root; EtymDB has no entry.
\end{cousinweb}

\begin{grammarlens}[title={What you've built}]
One more word for telling, explaining and describing at B1.
\end{grammarlens}

\subsection*{Your turn}
\begin{itemize}
\raggedright
  \item \emph{Say it:} \emph{pleskátsya}
  \item \emph{Say it:} \emph{pleskátsya}, once more
  \item \emph{Say it:} a sentence of your own with \emph{pleskátsya}
  \item \emph{Recall:} say the Russian for to slide, then the Russian for flippers, then say \emph{pleskátsya} again
\end{itemize}

\begin{tcolorbox}[breakable,colback=teal!4,colframe=teal!35!black,title={Before you move on}]
What does \ru{плескаться} mean? (\textquotedblleft{}To splash about\textquotedblright{}.) Say it once more.
\end{tcolorbox}
